% Chapter 5: The Check-In
% Scene function: Tert meets the Manager (Prim) for the first time.
% The bureaucratic review of the case. The gap between the form and
% the reality. Meri follow-up. Short, punchy, all office register.

\chapter{Good Acoustics}

My supervisor is designated Primus Ordinatus, which everyone shortens to Prim. It suits them. Prim is, technically, several hundred demons operating as a single administrative entity. This is not uncommon at the management level. It explains the thoroughness.

Prim stops by my desk on Tuesday morning. This means either there is a scheduling conflict or they have read one of my shift reports. Both would be unusual. Prim is carrying a clipboard. They always are. The mug says WORLD'S BEST and nothing else. It was issued to all middle management during a morale initiative whose purpose and origin nobody remembers.

``Tert,'' they say. ``Do you have a minute?''

I have exactly as many minutes as Prim needs me to have, because the question is rhetorical and the correct answer to a rhetorical question from your supervisor is always the one they already assume.

``Sure,'' I say.

``I wanted to touch base on forty-seven-C. The new atmospheric.''

``All right.''

``Intake flagged it as standard. Low religious engagement, caregiver background, single-occupant dwelling. How are you finding the environmental parameters?''

``Good acoustics,'' I say.

Prim writes this down. I watch them write it. They write ``good acoustics'' on a clipboard in a column I cannot see the heading of. There is a form for this conversation. This is not a casual check-in. This is a documented Week 1 Field Assessment. Someone in Compliance decided these needed to exist. There is now a paper trail for the fact that I said ``good acoustics'' to my supervisor on a Tuesday morning.

``And the Phase 1 benchmarks. Are we on track for the week-one Fear Index baseline?''

The Fear Index baseline for week one is a number pulled from a template that has no relationship to anything happening in the house. But it is on Prim's form, and Prim needs to fill in the field.

``On track,'' I say.

``Good. I see you modified the week-three escalation protocol. Direct interface felt premature?''

Someone read my modification. It turns out I preferred it when nobody did.

``For a caregiver subject with high baseline fatigue, yes.''

``That is good field judgment. Flag it for the mid-cycle so the reasoning is on record.''

``Done.''

Behind Prim, something in the corridor catches the light in a way that light should not be caught. Neither of us looks at it.

``One more thing. The Soul Outcome projection. Intake has this one at seventy-two percent positive. Any early reads that would adjust that?''

Seventy-two percent positive. For a secular subject with no church attendance, no religious framework, and no apparent relationship with God to sever. Seventy-two percent of what. Measured against what. Calibrated by the same methodology that produced the Fear Index of 7.2, which is calibrated against a meeting from 1987 that nobody has the minutes for. I could explain this to Prim. I have considered it. The explanation would be received, documented on the clipboard, filed, and would change nothing about the number on the form or the form's existence or Prim's belief in the form's utility.

``I will have more data by week two,'' I say.

``Perfect,'' Prim says. ``Keep me posted.''

They leave. The conversation took four minutes. Prim has a completed form, which is what Prim came for, and I am not entirely sure that is the small thing I used to think it was.

On my way to the transit desk I pass Meri in the hallway. She is not carrying the Operations Manual, which is either progress or surrender.

``I read the section,'' she says.

I do not need to ask which section. ``What did it say?''

``It said the mechanism is well-established and does not require field-level understanding.''

``There you go,'' I say.

``That is not an answer.''

``No,'' I say.
